\begin{lema}[Scaling after one move]\label{lem:hijo-fisico}
Consider a regular parent of size \(m\) whose physical point \(p^{Q}\) is positive on its binary support,
has unit row and column margins, satisfies \(p^{Q}_{ic}\le C/m\), and has incidence gap at least
\(\gamma>0\). Fix a legal cell \(e\), form the rook minor \(M_e\), and forbid any subset of the other
cells on the two diagonals of \(e\). For sufficiently large \(m\), all these masks of size \(m-1\) have
scalings in a common diffusion and gap class, with constants depending only on \(C,\gamma\). In
particular this holds for the full child \(M_e^{\rm fis}\) at first exit.
\end{lema}
\begin{proof}
We delete one column in each retained row and up to one cell per diagonal. The loss in each row is at
most \(3C/m\), and the same holds per column. Let \(R\) be the restriction of \(p^Q\) after the deletions
and let \(\mathcal W\) be the side with fewer vertices of a cut in the remaining bipartite graph. In the
parent, the Rayleigh bound on the centered indicator gives capacity at least \(\gamma |\mathcal W|/2\),
since \(|\mathcal W|\le m\). The capacity remaining in \(R\) is at least \((\gamma/2-3C/m)|\mathcal W|\).
Orient the edges from rows to columns. The row deficits and the negatives of the column deficits sum to
zero, and each has absolute value at most \(3C/m\). With \(\lambda=(3C/m)/(\gamma/2-3C/m)<1\), each cut
supports the net demand with symmetric capacities \(\lambda R\). The flow criterion gives a signed
correction and a doubly stochastic matrix \((1-\lambda)R\le p'\le(1+\lambda)R\). Its conductance is at
least \((1-\lambda)(\gamma/2-3C/m)\). Cheeger's inequality gives a uniformly positive gap for the
normalized Laplacian.

It remains to transfer these bounds from a feasible table to the scaling of the binary support. At size
\(k\) let \(a\) be doubly stochastic, positive on the support, with \(a_{ic}\le C'/k\) and gap
\(\gamma'\). The entropy maximum on this support is interior by strict feasibility. Its equations give
the scaling \(D\) with \(\log(kD_{ic})=u_i-u_c\). By the margins, \(\sum a_{ic}\log(kD_{ic})=\sum
D_{ic}\log(kD_{ic})\ge0\). The positive part of the first sum is at most \(\sum a_{ic}kD_{ic}\le C'k\).
Subtracting a median of the potentials and integrating over the level cuts gives
\[
 \frac{\gamma'}2\sum_z|u_z|\le\sum_{ic}a_{ic}|u_i-u_c|\le2C'k .
\]
Fix \(T>8(C')^2/\gamma'\). Each line has at least \(k/C'\) neighbors, and in the entire graph fewer than
\(k/(2C')\) vertices satisfy \(|u_z|>T\). The row equation \(1=e^{u_i}\sum_{c\sim i} e^{-u_c}/k\) gives
\(e^{u_i}\le2C'e^T\), and the column equation gives \(e^{-u_c}\le2C'e^T\). The same equations with the
upper bounds give reciprocal lower bounds. Thus \(d_-/k\le D_{ic}\le d_+/k\) on the support, and
\(D\ge(d_-/C')a\) preserves a positive gap by Rayleigh. Apply this to \(p'\) with \(k=m-1\). For the
parent mask it applies directly to \(p^Q\), including on regular long arrival. The predictor
margins~\cite{pc} hold up to the stopping time by Corollary~4.12 and the inner class of~\cite[Def.~2.2
and §§4--5]{tc}. If the gap is stated with response weights, their comparability with \(p^Q\) on the
positive-vacancy branch gives the Rayleigh hypothesis.
\end{proof}

The coefficient \(\coefT=\mathfrak F-1/12\) in (T) depends on the marked row and column masses \(s,t\)
and the crossing weight \(z\). The integrated expressions in Appendix~\ref{app:rook} are holomorphic and
bounded on disks with class radii. At \(s=0\) or \(t=0\) the visible terms are independent of \(z\).
Differentiating with respect to \(z\) makes both faces vanish. If \(M_0\) bounds \(\coefT\) on a polydisk
of class radius \(r\) about the deletion point in \(z\), Cauchy's estimate in \(z\) bounds
\(|\partial_z\coefT|\) by \(2M_0/r\) on the smaller disk. Schwarz's lemma in the masses \(s,t\) supplies
the factors \(|s|/r,|t|/r\). Hence \(\partial_z\coefT=st\,O(1)\) uniformly, including when integrating in
\(z\) along the deletion segment. Both masses equal \(1/k\). At size \(k\) this gives \(O(k^{-3})\) for
the change in \(\coefT/k\) per forbidden cell. Removing a row and a column changes the masses by
\(O(k^{-1})\) and the inverse size by \(O(k^{-2})\), at cost \(O(k^{-2})\) in \(\coefT/k\). There are at
most \(2m\) forbidden cells. The lemma provides the same margins at every step. Summing the changes and
using the remainder \(O(m^{-2})\) of (T) at both endpoints gives, for every child of the regular parent,
\begin{equation}\label{eq:incremento-hijo-fisico}
 |g_T(M_e^{\rm fis})-g_T(M)|\le C/m^2,\qquad |g_T(M_e^{\rm fis})|\le C/(m-1).
\end{equation}
